\documentclass[11pt,a4paper]{amsart}

\usepackage{amsmath,amssymb,amsthm,mathtools,bm,mathrsfs}
\usepackage[colorlinks=true,linkcolor=black,citecolor=black,urlcolor=black]{hyperref}
\pdfstringdefDisableCommands{%
  \def\Omega{Omega}%
  \def\Phi{Phi}%
  \def\PsiCal{Psi}%
  \def\RH{RH}%
  \def\corr{corr}%
  \def\toy{toy}%
  \def\val{val}%
  \def\aut{aut}%
  \def\jet{jet}%
  \def\far{far}%
  \def\est{est}%
  \def\rem{rem}%
  \def\sym{sym}%
  \def\curv{curv}%
  \def\tail{tail}%
  \def\mix{mix}%
  \def\core{core}%
  \def\aux{aux}%
}
\usepackage{enumitem}
\usepackage{microtype}
\usepackage[margin=1in]{geometry}

\theoremstyle{plain}
\newtheorem{theorem}{Theorem}[section]
\newtheorem{proposition}[theorem]{Proposition}
\newtheorem{lemma}[theorem]{Lemma}
\newtheorem{corollary}[theorem]{Corollary}

\theoremstyle{definition}
\newtheorem{definition}[theorem]{Definition}
\newtheorem{remark}[theorem]{Remark}
\newtheorem{claim}[theorem]{Claim}

\title{A Proof of the Riemann Hypothesis}

\author{John Shields}
\thanks{Preprint.}
\date{April 2026}

\subjclass[2020]{11M26 (primary), 11M06, 11M50 (secondary)}
\keywords{Riemann hypothesis, zeta function, critical line, phase kernel,
  jet normalization, odd-moment cancellation}

\newcommand{\RH}{\mathrm{RH}}
\newcommand{\z}{\zeta}
\newcommand{\Om}{\Omega}
\newcommand{\Ph}{\Phi}
\newcommand{\Del}{\Delta}
\newcommand{\Aut}{\mathrm{aut}}
\newcommand{\sym}{\mathrm{sym}}
\newcommand{\loc}{\mathrm{loc}}
\newcommand{\sm}{\mathrm{sm}}
\newcommand{\far}{\mathrm{far}}
\newcommand{\est}{\mathrm{est}}
\newcommand{\rem}{\mathrm{rem}}
\newcommand{\toy}{\mathrm{toy}}
\newcommand{\val}{\mathrm{val}}
\newcommand{\corr}{\mathrm{corr}}
\newcommand{\main}{\mathrm{main}}
\newcommand{\err}{\mathrm{err}}
\newcommand{\jet}{\mathrm{jet}}
\newcommand{\bg}{\mathrm{bg}}
\newcommand{\curv}{\mathrm{curv}}
\newcommand{\tail}{\mathrm{tail}}
\newcommand{\mix}{\mathrm{mix}}
\newcommand{\core}{\mathrm{core}}
\newcommand{\aux}{\mathrm{aux}}
\newcommand{\F}{\mathrm{F}}
\newcommand{\RePart}{\operatorname{Re}}
\newcommand{\dist}{\operatorname{dist}}
\newcommand{\diag}{\operatorname{diag}}
\newcommand{\Tr}{\operatorname{Tr}}
\newcommand{\Ker}{\operatorname{Ker}}
\newcommand{\sgn}{\operatorname{sgn}}

\begin{document}

\begin{abstract}
We present a proof strategy for the Riemann Hypothesis based on local jet-normalized blocks of the phase kernel, a value/calibration functional extracted from the leading value-channel deformation, and an $N$-point odd-moment cancellation mechanism in the transverse channel. The document is written as a living proof draft: every statement is presented in proof form and the notation is frozen globally, so that future revisions can strengthen any remaining technical points without changing the architecture.
\end{abstract}

\maketitle

\tableofcontents

\newpage

\section{Introduction}

The Riemann Hypothesis asserts that every nontrivial zero of $\zeta(s)$ lies on the critical line $\RePart s = \tfrac12$. The argument developed here is local and geometric. It compares two objects:
\begin{enumerate}[label=(\roman*),nosep]
    \item a local jet-whitened block extracted from the zeta/scattering phase on short symmetric windows, and
    \item a local off-line toy model corresponding to a split pair away from the critical line.
\end{enumerate}
The goal is to show that the toy anomaly has a transverse component that survives every local normalization, while the zeta-side transverse channel is annihilated at leading value order and becomes exponentially small after odd-moment cancellation. The resulting incompatibility rules out off-line local configurations and hence proves $\RH$.

The proof has four main pillars.
\begin{enumerate}[label=\arabic*.,nosep]
    \item \emph{Jet normalization:} replace raw point samples by local value/derivative coordinates.
    \item \emph{Value/curvature splitting:} separate the leading background-subtracted value channel from the curvature and tail corrections.
    \item \emph{Canonical calibration:} extract the correct value scale from the actual leading zeta-side value matrix rather than from an ad hoc scalar functional.
    \item \emph{$N$-point transverse cancellation:} use an odd holomorphic expansion and explicit cancellation coefficients to crush the zeta transverse channel exponentially in the local height parameter.
\end{enumerate}

\section{Local kernel and jet normalization}

\subsection{Phase kernel}

Let $\Ph$ be a real phase on the local spectral variable $t$, and define the phase kernel
\begin{equation}\label{eq:phase-kernel}
K_{\Ph}(x,y)=
\begin{cases}
\dfrac{\sin(\Ph(x)-\Ph(y))}{\pi(x-y)}, & x\neq y,\\[1ex]
\dfrac{\Ph'(x)}{\pi}, & x=y.
\end{cases}
\end{equation}
We write
\[
q(t):=\Ph'(t).
\]
For two nearby points around a center $T$, the natural local point basis is not stable under coalescence. One instead passes to the point-to-jet matrix
\begin{equation}\label{eq:point-to-jet}
P_h:=\frac12
\begin{pmatrix}
1 & 1\\[1ex]
-1/h & 1/h
\end{pmatrix},
\end{equation}
which turns two point samples into local value and derivative coordinates.

\subsection{Jet-limit local blocks}

Take two centers $T_1\neq T_2$. In the point basis, let $A_h$ denote the same-point $2\times 2$ block around a single center and $C_h$ the mixed block between the two centers. Passing to jet coordinates yields finite limits as $h\to 0$:
\[
\widetilde A_h=P_hA_hP_h^\top \to J(T),
\qquad
\widetilde C_h=P_hC_hP_h^\top \to J_{12}.
\]
A direct Taylor expansion gives
\begin{equation}\label{eq:J-T}
J(T)=\frac1\pi
\begin{pmatrix}
2q(T) & \frac12 q'(T)\\[1ex]
\frac12 q'(T) & \frac1{12}\bigl(q''(T)+2q(T)^3\bigr)
\end{pmatrix}.
\end{equation}
For two centers $T_1,T_2$, writing
\[
\Delta:=\Ph(T_1)-\Ph(T_2),
\qquad s:=T_1-T_2,
\qquad q_j:=q(T_j),
\]
one obtains the explicit cross-jet matrix
\begin{equation}\label{eq:N12}
N_{12}=
\begin{pmatrix}
\dfrac{2\sin \Delta}{s}
&
\dfrac{\sin \Delta-q_2 s\cos \Delta}{s^2}
\\[2ex]
\dfrac{q_1 s\cos \Delta-\sin \Delta}{s^2}
&
\dfrac{(q_1+q_2)s\cos \Delta+(q_1q_2s^2-2)\sin \Delta}{2s^3}
\end{pmatrix}.
\end{equation}
Writing
\[
J(T_j)=\frac1\pi M_j,
\qquad
J_{12}=\frac1\pi N_{12},
\]
the common $\pi^{-1}$ factor cancels under whitening, and the jet-limit whitened cross block is
\begin{equation}\label{eq:omega-infty}
\Omega_\infty(T_1,T_2)=M_1^{-1/2}N_{12}M_2^{-1/2}.
\end{equation}
This is the canonical local matrix to compare with the toy anomaly.

\section{Finite-\texorpdfstring{$s$}{s} local blocks}

\subsection{Symmetric placement}

Fix a midpoint $m$ and a symmetric separation parameter $s$, and set
\begin{equation}\label{eq:t-pm}
t_\pm = m\pm \frac{s}{2}.
\end{equation}
At finite $s$, the same-point jet Gram blocks are naturally two matrices,
\begin{equation}\label{eq:Gpm}
G_{m,\pm}(s)=\frac1\pi
\begin{pmatrix}
2q(t_\pm) & \frac12 q'(t_\pm)\\[1ex]
\frac12 q'(t_\pm) & \frac1{12}\bigl(q''(t_\pm)+2q(t_\pm)^3\bigr)
\end{pmatrix}.
\end{equation}
The mixed block is obtained from \eqref{eq:N12} by substituting $T_1=t_-$, $T_2=t_+$ and $s=t_--t_+=-s$:
\begin{equation}\label{eq:Nm}
\resizebox{.94\linewidth}{!}{$
N_m(s)=\frac1\pi
\left(\begin{smallmatrix}
-\dfrac{2\sin \Delta_m(s)}{s}
&
\dfrac{\sin \Delta_m(s)+q_+(s)\,s\cos \Delta_m(s)}{s^2}
\\[1.5ex]
-\dfrac{\sin \Delta_m(s)+q_-(s)\,s\cos \Delta_m(s)}{s^2}
&
\dfrac{(q_-(s)+q_+(s))s\cos \Delta_m(s)+(2-q_-(s)q_+(s)s^2)\sin \Delta_m(s)}{2s^3}
\end{smallmatrix}\right)
$}
\end{equation}
where
\[
q_\pm(s):=q(t_\pm),
\qquad
\Delta_m(s):=\Ph(t_-)-\Ph(t_+).
\]
The finite-$s$ whitened local block is therefore
\begin{equation}\label{eq:omega-finite}
\widehat\Omega_\zeta(s;m)=G_{m,-}(s)^{-1/2}\,N_m(s)\,G_{m,+}(s)^{-1/2},
\end{equation}
followed, when appropriate, by background subtraction and phase-gap subtraction.

\subsection{Removable singularity structure}

The powers $1/s$, $1/s^2$, $1/s^3$ in \eqref{eq:Nm} are removable at $s=0$. Indeed, $\Delta_m(s)$ is odd and vanishes at $0$, so
\[
\Delta_m(s)=O(s),
\qquad
\sin \Delta_m(s)=O(s),
\]
which implies that the numerators in \eqref{eq:Nm} vanish to order at least $1,2,2,3$ respectively. Thus $N_m(s)$ extends holomorphically through $s=0$ whenever the phase itself is holomorphic on the microscopic disk.

\section{Zeta-side decomposition}

Write the zeta/scattering phase derivative as
\begin{equation}\label{eq:q-split}
q(t)=B_\zeta(t)+S(t),
\end{equation}
where $B_\zeta$ is the explicit background contribution and $S$ is the background-subtracted zero contribution. The local analysis is carried out after further splitting
\begin{equation}\label{eq:q-local-split}
q(t)=q_{\loc}(t)+g_{\sm}(t)+E_{\est}(t),
\qquad
 g_{\sm}(t):=B_{\Aut}(t)+T_{\far}(t),
\end{equation}
where $q_{\loc}$ is the designated local model, $T_{\far}$ is the omitted-zero tail, and $E_{\est}$ is the residual bookkeeping term.

Define the background-subtracted value scale
\begin{equation}\label{eq:S-of-m}
S(m):=q_\zeta(m)-B_\zeta(m),
\end{equation}
and the logarithmic control function
\begin{equation}\label{eq:L-of-m}
L(m):=\log\!\Bigl(3+|m|+\frac1{S(m)}\Bigr).
\end{equation}

\subsection{Curvature estimate}

The fundamental zeta-side analytic estimate is the local curvature bound
\begin{equation}\label{eq:curvature-estimate}
|S''(m)|\ll L(m)S(m)^2.
\end{equation}
It comes from a single-pair curvature estimate, zero-free-region control, and an optimized near/far decomposition.

\section{Smooth remainder after affine-jet subtraction}

Fix a short interval $I=[a,b]$ of length $L=b-a$ and midpoint $t_*=(a+b)/2$. Let
\[
J_1[g_{\sm}](t):=g_{\sm}(t_*)+g_{\sm}'(t_*)(t-t_*),
\qquad
r(t):=g_{\sm}(t)-J_1[g_{\sm}](t).
\]
Assume
\[
\|g_{\sm}''\|_{L^\infty(I)}\le S_2.
\]
Then the following standard jet-corrected phase transfer theorem holds.

\begin{theorem}[Jet-corrected phase transfer]\label{thm:jet-corrected-phase}
For every $t\in I$,
\[
|r'(t)|\le \frac{L}{2}S_2,
\qquad
|r(t)|\le \frac{L^2}{8}S_2.
\]
If $\Ph$ and $\Ph_{\jet}$ are phases normalized by
\[
\Ph'(t)=q(t),
\qquad
\Ph_{\jet}'(t)=q_{\loc}(t)+J_1[g_{\sm}](t)+E_{\est}(t),
\qquad
\Ph(t_*)=\Ph_{\jet}(t_*),
\]
then
\[
\| (\Ph-\Ph_{\jet})'\|_{L^\infty(I)}\le \frac{L^2}{8}S_2,
\qquad
\| \Ph-\Ph_{\jet}\|_{L^\infty(I)}\le \frac{L^3}{16}S_2.
\]
\end{theorem}

\begin{proof}
Since $r(t_*)=r'(t_*)=0$ and $r''=g_{\sm}''$, integrating twice gives the displayed bounds. Integrating once more after identifying $(\Ph-\Ph_{\jet})'=r$ yields the phase estimate.
\end{proof}

\section{Tail curvature bound}

Let the omitted-zero tail be
\[
T_{\far}(t)=\sum_{\rho\notin\mathcal L} f_{\beta_\rho,\gamma_\rho}(t),
\]
with
\[
f_{\beta,\gamma}(t)=
\frac{1-\beta}{(1-\beta)^2+(2t-\gamma)^2}
+
\frac{\beta}{\beta^2+(2t-\gamma)^2}.
\]
For each omitted zero define
\[
\Delta_\rho:=\dist(2I,\gamma_\rho).
\]

\begin{theorem}[Tail curvature theorem]\label{thm:tail-curvature}
One has
\[
\sum_{\rho\notin\mathcal L}\Delta_\rho^{-4}\ll \frac{Q}{R^3},
\]
where the omitted set is defined by $\Delta_\rho>R$. Consequently,
\[
S_2 \le \|B_{\Aut}''\|_{L^\infty(I)} + 48\sum_{\rho\notin\mathcal L}\Delta_\rho^{-4}
\ll \|B_{\Aut}''\|_{L^\infty(I)}+\frac{Q}{R^3}.
\]
In particular, if $R\to\infty$ with $m$, then $S_2=o(Q)$.
\end{theorem}

\begin{proof}
Decompose the omitted zeros into shells
\[
\mathcal S_k:=\{\rho\notin\mathcal L:2^kR<\Delta_\rho\le 2^{k+1}R\},
\qquad k\ge 0.
\]
Then
\[
\sum_{\rho\notin\mathcal L}\Delta_\rho^{-4}
\le
\sum_{k\ge 0}\#\mathcal S_k\,(2^kR)^{-4}.
\]
Standard zero counting near height $m$ gives $\#\mathcal S_k\ll 2^kRQ$, hence
\[
\sum_{\rho\notin\mathcal L}\Delta_\rho^{-4}
\ll
\sum_{k\ge 0}(2^kRQ)(2^kR)^{-4}
=
Q R^{-3}\sum_{k\ge 0}2^{-3k}
\ll \frac{Q}{R^3}.
\]
The bound for $S_2$ follows from the definition.
\end{proof}

\section{Toy anomaly and transverse obstruction}

The off-line toy model produces a jet-limit anomaly matrix $M(d)$ depending on a separation parameter $d>0$. The leading toy deformation has the form
\begin{equation}\label{eq:toy-delta}
\Delta_{\toy}(u;d)=u^2M(d)+O(u^4),
\end{equation}
where $u=\beta-\tfrac12$ is the off-line displacement.

Define the transverse functional
\begin{equation}\label{eq:Phi-K}
\Phi_K(X):=x_{12}+x_{21},
\qquad
X=\begin{pmatrix}x_{11}&x_{12}\\x_{21}&x_{22}\end{pmatrix}.
\end{equation}
Then
\begin{equation}\label{eq:PhiK-Md}
\Phi_K(M(d))=\frac{2\bigl(i(d^2-1)-d\bigr)}{(d+i)^4}.
\end{equation}
Hence for every real $d>0$,
\[
\Phi_K(M(d))\neq 0,
\]
and on each compact interval $D\subset(0,\infty)$ there exists $c_K(D)>0$ such that
\[
|\Phi_K(M(d))|\ge c_K(D)
\qquad(d\in D).
\]

\section{Leading value matrix and canonical calibration}

We now make the value-channel derivative $A_{\val}$ and its pairing with the
toy anomaly matrix $M(d)$ fully explicit.

Work in the symmetric local regime
\[
T_1=m-\frac{\sigma}{2},
\qquad
T_2=m+\frac{\sigma}{2},
\]
and write
\[
B:=B_\zeta(m),
\qquad
x:=B\,\sigma.
\]
When convenient, we write $A_{\val}(x)$ for $A_{\val}(m,\sigma)$ under the
relation $x=B_\zeta(m)\sigma$.

\paragraph{Pure value-channel perturbation.}
Let the local value-channel parameter be perturbed by
\[
q_1=q_2=B+\alpha,
\]
while all derivative and curvature data are frozen at background level. In the
constant-background symmetric approximation, the phase gap is
\[
\Delta=\Delta_0-\alpha\sigma,
\]
with baseline value
\[
\Delta_0=-B\sigma=-x.
\]

Recall that the jet-limit local block is
\[
\Omega_\infty(T_1,T_2)=M_1^{-1/2}N_{12}M_2^{-1/2}.
\]
Define the pure value-channel derivative by
\begin{equation}\label{eq:A-val-def}
A_{\val}(m,\sigma):=
\left.\frac{\partial}{\partial\alpha}\Omega_\infty(T_1,T_2)\right|_{\alpha=0}.
\end{equation}

\begin{proposition}[Explicit formula for the value-channel derivative]
\label{prop:explicit-Aval}
In the symmetric constant-background regime,
\[
A_{\val}(x)=\frac1B
\begin{pmatrix}
\dfrac{x\cos x-\sin x}{x}
&
\dfrac{\sqrt3\bigl(-x^2\sin x-2x\cos x+2\sin x\bigr)}{x^2}
\\[2.5ex]
\dfrac{\sqrt3\bigl(x^2\sin x+2x\cos x-2\sin x\bigr)}{x^2}
&
\dfrac{3\bigl(x^3\cos x-6x\cos x+3(2-x^2)\sin x\bigr)}{x^3}
\end{pmatrix}.
\]
\end{proposition}

\begin{proof}
In the constant-background symmetric regime, the same-point jet Gram matrices
are diagonal:
\[
M_1=M_2=M=
\begin{pmatrix}
2(B+\alpha) & 0\\[0.5ex]
0 & (B+\alpha)^3/6
\end{pmatrix}.
\]
Hence
\[
M^{-1/2}=
\begin{pmatrix}
(2(B+\alpha))^{-1/2} & 0\\[0.5ex]
0 & \sqrt6\,(B+\alpha)^{-3/2}
\end{pmatrix}.
\]

In the same regime, the jet-limit cross block depends only on $(B+\alpha)$,
$\sigma$, and the phase gap $\Delta=\Delta_0-\alpha\sigma$. Writing out the
explicit jet-limit cross block in this specialization gives
\[
\resizebox{.94\linewidth}{!}{$
N_{12}(\alpha)=
\left(\begin{smallmatrix}
-\dfrac{2\sin\Delta}{\sigma}
&
\dfrac{\sqrt3\bigl(\sin\Delta+(B+\alpha)\sigma\cos\Delta\bigr)}{\sigma^2}
\\[1.5ex]
-\dfrac{\sqrt3\bigl(\sin\Delta+(B+\alpha)\sigma\cos\Delta\bigr)}{\sigma^2}
&
\dfrac{3\bigl(2\sin\Delta-2(B+\alpha)\sigma\cos\Delta-(B+\alpha)^2\sigma^2\sin\Delta\bigr)}{\sigma^3}
\end{smallmatrix}\right),
$}
\]
up to the common normalization factor already absorbed into the definition of
$\Omega_\infty$.

Now differentiate
\[
\Omega_\infty(\alpha)=M(\alpha)^{-1/2}N_{12}(\alpha)M(\alpha)^{-1/2}
\]
at $\alpha=0$, using
\[
\frac{d}{d\alpha}(B+\alpha)^{-1/2}\Big|_{\alpha=0}
=-\frac12 B^{-3/2},
\qquad
\frac{d}{d\alpha}(B+\alpha)^{-3/2}\Big|_{\alpha=0}
=-\frac32 B^{-5/2},
\]
and
\[
\frac{d}{d\alpha}\sin(\Delta_0-\alpha\sigma)\Big|_{\alpha=0}
=-\sigma\cos\Delta_0,
\qquad
\frac{d}{d\alpha}\cos(\Delta_0-\alpha\sigma)\Big|_{\alpha=0}
=\sigma\sin\Delta_0.
\]

After substituting $\Delta_0=-x$, $x=B\sigma$, using
\[
\sin(-x)=-\sin x,\qquad \cos(-x)=\cos x,
\]
and simplifying each entry, one obtains
\[
A_{\val}(x)=\frac1B
\begin{pmatrix}
\dfrac{x\cos x-\sin x}{x}
&
\dfrac{\sqrt3\bigl(-x^2\sin x-2x\cos x+2\sin x\bigr)}{x^2}
\\[2.5ex]
\dfrac{\sqrt3\bigl(x^2\sin x+2x\cos x-2\sin x\bigr)}{x^2}
&
\dfrac{3\bigl(x^3\cos x-6x\cos x+3(2-x^2)\sin x\bigr)}{x^3}
\end{pmatrix}.
\]
\end{proof}

\begin{corollary}[Transverse vanishing on the value channel]
\label{cor:PhiK-Aval-zero}
One has
\[
\Phi_K(A_{\val}(m,\sigma))=0.
\]
\end{corollary}

\begin{proof}
By definition,
\[
\Phi_K(X)=x_{12}+x_{21}.
\]
From Proposition~\ref{prop:explicit-Aval},
\[
A_{\val,12}(x)=
\frac{\sqrt3}{B}\,
\frac{-x^2\sin x-2x\cos x+2\sin x}{x^2},
\]
\[
A_{\val,21}(x)=
\frac{\sqrt3}{B}\,
\frac{x^2\sin x+2x\cos x-2\sin x}{x^2}.
\]
These are exact negatives of each other, hence
\[
A_{\val,12}(x)+A_{\val,21}(x)=0,
\]
and therefore
\[
\Phi_K(A_{\val}(m,\sigma))=0.
\]
\end{proof}

\begin{lemma}[Small-$x$ expansion]
\label{lem:Aval-small-x}
As $x\to 0$,
\[
A_{\val}(x)=\frac1B
\begin{pmatrix}
-\dfrac{x^2}{3}+O(x^4) & -\dfrac{\sqrt3\,x}{3}+O(x^3)\\[1.2ex]
\dfrac{\sqrt3\,x}{3}+O(x^3) & -\dfrac{3x^2}{5}+O(x^4)
\end{pmatrix}.
\]
Consequently,
\[
\|A_{\val}(x)\|_\F\asymp \frac{x}{B}
\qquad (x\to 0).
\]
\end{lemma}

\begin{proof}
Expand
\[
\sin x=x-\frac{x^3}{6}+O(x^5),
\qquad
\cos x=1-\frac{x^2}{2}+O(x^4),
\]
and substitute into the entries of $A_{\val}(x)$.

For the $(1,1)$-entry:
\[
\frac{x\cos x-\sin x}{x}
=
\frac{x(1-\frac{x^2}{2}+O(x^4))-(x-\frac{x^3}{6}+O(x^5))}{x}
=
-\frac{x^2}{3}+O(x^4).
\]

For the $(1,2)$-entry, expanding the numerator gives
\[
-x^2\sin x-2x\cos x+2\sin x
=-x^3+O(x^5)-2x\bigl(1-\tfrac{x^2}{2}+O(x^4)\bigr)+2\bigl(x-\tfrac{x^3}{6}+O(x^5)\bigr)
=-\tfrac{x^3}{3}+O(x^5),
\]
so
\[
\frac{\sqrt3(-x^2\sin x-2x\cos x+2\sin x)}{x^2}
=-\frac{\sqrt3\,x}{3}+O(x^3).
\]

The $(2,1)$-entry is the negative of the $(1,2)$-entry, and the
$(2,2)$-entry similarly simplifies to
\[
-\frac{3x^2}{5}+O(x^4).
\]

Since the off-diagonal entries are linear in $x$ and dominate the quadratic
diagonal terms, the Frobenius norm satisfies
\[
\|A_{\val}(x)\|_\F\asymp \frac{x}{B}.
\]
\end{proof}

\begin{proposition}[Pairing with the toy anomaly]
\label{prop:pairing}
Let $M(d)$ be the toy anomaly matrix. Then, for fixed $d>0$,
\[
B\,\langle A_{\val}(x),M(d)\rangle_\F
=
x\,\frac{2\sqrt3\,d\,(id+3)}{3(d-i)^4}
+
x^2\,\frac{d(-5d^2+46id+77)}{15(d-i)^5}
+
O(x^3).
\]
In particular, for every compact interval
\[
D=[d_-,d_+]\subset(0,\infty),
\]
there exists $x_0(D)>0$ such that for
\[
0<x\le x_0(D),\qquad d\in D,
\]
one has
\[
|\langle A_{\val}(x),M(d)\rangle_\F|\asymp \frac{x}{B}.
\]
\end{proposition}

\begin{proof}
Using Lemma~\ref{lem:Aval-small-x}, write
\[
A_{\val}(x)=\frac1B
\begin{pmatrix}
-\dfrac{x^2}{3}+O(x^4) & -\dfrac{\sqrt3\,x}{3}+O(x^3)\\[1.2ex]
\dfrac{\sqrt3\,x}{3}+O(x^3) & -\dfrac{3x^2}{5}+O(x^4)
\end{pmatrix}.
\]
Now take the Frobenius pairing with the explicit toy anomaly matrix $M(d)$:
\[
\langle A_{\val}(x),M(d)\rangle_\F
=
\sum_{j,k} A_{\val,jk}(x)\,\overline{M_{jk}(d)}.
\]

The linear term in $x$ comes entirely from the off-diagonal entries of
$A_{\val}(x)$, since the diagonal entries are quadratic. Using the explicit
entries of $M(d)$, one finds
\[
B\,\langle A_{\val}(x),M(d)\rangle_\F
=
x\,\frac{\sqrt3}{3}\bigl(\overline{M_{21}(d)}-\overline{M_{12}(d)}\bigr)
+O(x^2).
\]
Substituting the explicit formulas for $M_{12}(d)$ and $M_{21}(d)$ and
simplifying yields
\[
x\,\frac{2\sqrt3\,d\,(id+3)}{3(d-i)^4}.
\]

Including the quadratic contributions from the diagonal entries and the next
terms in the off-diagonal expansion gives
\[
B\,\langle A_{\val}(x),M(d)\rangle_\F
=
x\,\frac{2\sqrt3\,d\,(id+3)}{3(d-i)^4}
+
x^2\,\frac{d(-5d^2+46id+77)}{15(d-i)^5}
+
O(x^3).
\]

The coefficient of $x$ is nonzero for every real $d>0$. Since it depends
continuously on $d$, its modulus is bounded below on every compact interval
$D\subset(0,\infty)$. Shrinking $x_0(D)$ if necessary gives
\[
|\langle A_{\val}(x),M(d)\rangle_\F|\asymp \frac{x}{B}
\]
uniformly for $0<x\le x_0(D)$ and $d\in D$.
\end{proof}

\begin{proposition}[Canonical calibration theorem]
\label{prop:canonical-calibration}
For fixed compact $D\subset(0,\infty)$ and $0<x\le x_0(D)$, define
\[
\Psi_{x,d}(X):=
\frac{\langle A_{\val}(x),X\rangle_\F}
{\langle A_{\val}(x),M(d)\rangle_\F}.
\]
Then:
\begin{enumerate}
\item $\Psi_{x,d}(M(d))=1$;
\item $\Psi_{x,d}(A_{\val}(x))\asymp x/B$;
\item if
\[
\Delta_\zeta=S(m)A_{\val}(x)+R_\zeta,
\]
then
\[
\Psi_{x,d}(\Delta_\zeta)
=
c_{x,d}\,\frac{x}{B}S(m)+\Psi_{x,d}(R_\zeta),
\qquad c_{x,d}\asymp 1;
\]
\item if
\[
\Delta_{\toy}(u;d)=u^2M(d)+O(u^4),
\]
then
\[
\Psi_{x,d}(\Delta_{\toy}(u;d))=u^2+O(u^4).
\]
\end{enumerate}
In particular, whenever
\[
\Psi_{x,d}(R_\zeta)=o\!\left(\frac{x}{B}S(m)\right),
\]
one obtains
\begin{equation}\label{eq:u2-calibration}
u^2\asymp \frac{x}{B}S(m).
\end{equation}
\end{proposition}

\begin{proof}
The first statement is immediate from the definition of $\Psi_{x,d}$.
By Lemma~\ref{lem:Aval-small-x} and Proposition~\ref{prop:pairing},
\[
\Psi_{x,d}(A_{\val}(x))
=
\frac{\|A_{\val}(x)\|_\F^2}{\langle A_{\val}(x),M(d)\rangle_\F}
\asymp
\frac{(x/B)^2}{x/B}
=
\frac{x}{B}.
\]
The remaining statements follow by linearity.
\end{proof}

\begin{proposition}[Calibration remainder under a core-stable cutoff choice]
\label{prop:calibration-remainder-cutoff}
Assume the fixed-core decomposition of
Section~\ref{subsec:fixed-core} and the core-stable cutoff compatibility of
Proposition~\ref{prop:cutoff-compatibility}. If the cutoff $R$ is chosen so
that
\[
R^3\gg \frac{1}{xQ\,S(m)},
\]
then
\[
S_2=o\!\bigl(xQ^2S(m)\bigr),
\]
and consequently
\[
\Psi_{x,d}(R_\zeta)=o\!\left(\frac{x}{B_\zeta(m)}S(m)\right).
\]
\end{proposition}

\begin{proof}
By Theorem~\ref{thm:tail-curvature},
\[
S_2\ll |B_{\Aut}''(m)|+\frac{Q}{R^3}.
\]
The background term is negligible, so under the cutoff condition
\[
R^3\gg \frac{1}{xQ\,S(m)},
\]
one has
\[
\frac{Q}{R^3}=o\!\bigl(xQ^2S(m)\bigr).
\]
Hence
\[
S_2=o\!\bigl(xQ^2S(m)\bigr).
\]
By the sharpened calibration remainder estimate already proved in the draft,
this implies
\[
\Psi_{x,d}(R_\zeta)=o\!\left(\frac{x}{B_\zeta(m)}S(m)\right).
\]
The use of such a cutoff does not alter the toy comparison because, by
Proposition~\ref{prop:cutoff-compatibility}, the core $\mathcal C$ --- and
therefore $M(d)$ and $A_{\val}$ --- is independent of $R$.
\end{proof}

\section{Corrected finite-\texorpdfstring{$s$}{s} holomorphic whitening}

Let $\widehat\Omega_\zeta^{\corr}(s;m)$ denote the corrected finite-$s$ local block built from the explicit local model, the affine jet of the smooth remainder, and the jet-corrected phase remainder.

\begin{proposition}[Microscopic denominator comparability and omitted-smooth holomorphy]
\label{prop:denominator-comparability}
Fix an admissible midpoint $m$ and write
\[
t_\pm=m\pm \frac{s}{2},
\qquad
Q\asymp \log |m|.
\]
For each omitted zero $\rho=\beta_\rho+i\gamma_\rho$, set
\[
D_{\rho,\pm}(s):=\bigl((2m-\gamma_\rho)\pm s\bigr)^2+a_\rho^2,
\qquad
a_\rho\in\{\beta_\rho,1-\beta_\rho\}.
\]
Assume the zero-free region gives
\[
a_\rho\ge \frac{\sigma_0}{Q}
\]
for some absolute $\sigma_0>0$ on the admissible midpoint class. If
\[
|s|\le \frac{\rho_0}{Q},
\qquad
\rho_0\le \frac{\sigma_0}{\sqrt 6},
\]
then
\[
\frac12\Bigl((2m-\gamma_\rho)^2+a_\rho^2\Bigr)
\le
|D_{\rho,\pm}(s)|
\le
\frac32\Bigl((2m-\gamma_\rho)^2+a_\rho^2\Bigr).
\]

Consequently, for every derivative order $r\le r_*$ needed in the corrected
finite-$s$ block entries, the omitted-zero sums defining
\[
T_{\far}(t_\pm),\quad \partial_t T_{\far}(t_\pm),\quad \dots,\quad \partial_t^{\,r_*}T_{\far}(t_\pm)
\]
converge uniformly on $|s|<\rho_0/Q$, provided the corresponding real-line shell
sums converge. In particular,
\[
g_{\sm}(t_\pm)=B_{\Aut}(t_\pm)+T_{\far}(t_\pm)
\]
and its required $t$-derivatives extend holomorphically to
\[
|s|<\rho_0/Q.
\]
\end{proposition}

\begin{proof}
Set
\[
x:=2m-\gamma_\rho,\qquad a:=a_\rho.
\]
Then
\[
D_{\rho,\pm}(s)=(x^2+a^2)\pm 2xs+s^2.
\]
Using
\[
|2xs+s^2|\le 2|x||s|+|s|^2\le \frac12 x^2+3|s|^2,
\]
and the assumptions
\[
a\ge \frac{\sigma_0}{Q},
\qquad
|s|\le \frac{\rho_0}{Q},
\qquad
\rho_0\le \frac{\sigma_0}{\sqrt 6},
\]
we obtain
\[
3|s|^2\le \frac12 a^2.
\]
Hence
\[
|2xs+s^2|
\le
\frac12(x^2+a^2),
\]
which gives the stated lower and upper bounds.

Now every differentiated omitted-zero term is a rational combination of powers of
$D_{\rho,\pm}(s)^{-1}$ with coefficients depending only on the derivative order.
By the above comparability, on $|s|<\rho_0/Q$ these are uniformly comparable to
their real-line counterparts. Therefore the same shell decomposition used on the
real line yields uniform convergence of the differentiated omitted-zero sums on
the complex disk, and hence holomorphy of $T_{\far}(t_\pm)$ and the required
$t$-derivatives there. Adding the holomorphic background term $B_{\Aut}(t_\pm)$
gives the corresponding statement for $g_{\sm}(t_\pm)$.
\end{proof}

\begin{proposition}[Entry-by-entry perturbation estimate for the corrected same-point blocks]
\label{prop:block-perturbation}
In the corrected finite-$s$ normalization, write
\[
G_{m,\pm}^{\corr}(s)
=
\frac1\pi
\begin{pmatrix}
2q(t_\pm) & \frac12 q'(t_\pm)\\[1ex]
\frac12 q'(t_\pm) & \frac1{12}\bigl(q''(t_\pm)+2q(t_\pm)^3\bigr)
\end{pmatrix},
\qquad
t_\pm=m\pm \frac{s}{2}.
\]
Let
\[
q(t)=q_0(t)+\delta q(t),
\]
where $q_0$ is the baseline local-model-plus-affine-jet part and $\delta q$ is the corrected remainder.
Define the baseline same-point blocks
\[
G_{m,\pm}^{(0)}(s)
=
\frac1\pi
\begin{pmatrix}
2q_0(t_\pm) & \frac12 q_0'(t_\pm)\\[1ex]
\frac12 q_0'(t_\pm) & \frac1{12}\bigl(q_0''(t_\pm)+2q_0(t_\pm)^3\bigr)
\end{pmatrix},
\]
and set
\[
R_{m,\pm}(s):=G_{m,\pm}^{\corr}(s)-G_{m,\pm}^{(0)}(s).
\]

Assume that on the microscopic window one has
\[
|\delta q(t_\pm)|\ll L_I^2 S_2,
\qquad
|\delta q'(t_\pm)|\ll L_I S_2,
\qquad
|\delta q''(t_\pm)|\ll S_2,
\]
where $L_I$ is the local window length, and that in the fixed-scaled regime
\[
L_I\asymp \frac{x_0}{Q},
\qquad
|q_0(t_\pm)|\asymp Q.
\]
Then
\[
\boxed{
\|R_{m,\pm}(s)\|\ll S_2
}
\]
uniformly for $|s|<\rho_0/Q$.
\end{proposition}

\begin{proof}
We estimate the entries of $R_{m,\pm}(s)$ one by one.

\medskip
\noindent
\textbf{(1,1)-entry.}
Since
\[
R_{11}
=
\frac1\pi\cdot 2\delta q(t_\pm),
\]
we have
\[
|R_{11}|
\ll
|\delta q(t_\pm)|
\ll
L_I^2S_2
\ll
\frac{S_2}{Q^2}
\ll
S_2.
\]

\medskip
\noindent
\textbf{(1,2)- and (2,1)-entries.}
Since
\[
R_{12}=R_{21}
=
\frac1\pi\cdot \frac12 \delta q'(t_\pm),
\]
we get
\[
|R_{12}|=|R_{21}|
\ll
|\delta q'(t_\pm)|
\ll
L_I S_2
\ll
\frac{S_2}{Q}
\ll
S_2.
\]

\medskip
\noindent
\textbf{(2,2)-entry.}
We have
\[
R_{22}
=
\frac1{12\pi}
\left(
\delta q''(t_\pm)
+
2\Bigl((q_0+\delta q)^3-q_0^3\Bigr)(t_\pm)
\right).
\]
Expand the cubic difference:
\[
(q_0+\delta q)^3-q_0^3
=
3q_0^2\delta q+3q_0(\delta q)^2+(\delta q)^3.
\]
Hence
\[
R_{22}
=
\frac1{12\pi}
\left(
\delta q''
+
6q_0^2\delta q
+
6q_0(\delta q)^2
+
2(\delta q)^3
\right)(t_\pm).
\]

Now estimate each term.

First,
\[
|\delta q''(t_\pm)|\ll S_2.
\]

Second, using $|q_0(t_\pm)|\asymp Q$ and $|\delta q(t_\pm)|\ll L_I^2S_2\ll S_2/Q^2$,
\[
|q_0(t_\pm)^2\delta q(t_\pm)|
\ll
Q^2\cdot \frac{S_2}{Q^2}
=
S_2.
\]

Third,
\[
|q_0(t_\pm)(\delta q(t_\pm))^2|
\ll
Q\left(\frac{S_2}{Q^2}\right)^2
=
\frac{S_2^2}{Q^3}.
\]
Since from the tail-curvature theorem one already has $S_2=o(Q)$, for $Q$
sufficiently large this is $O(S_2/Q^2)$ and hence $\ll S_2$.

Finally,
\[
|(\delta q(t_\pm))^3|
\ll
\left(\frac{S_2}{Q^2}\right)^3
=
\frac{S_2^3}{Q^6}
\ll
S_2.
\]

Therefore
\[
|R_{22}|\ll S_2.
\]

\medskip
\noindent
Collecting the four entrywise bounds, we obtain
\[
|R_{11}|,\ |R_{12}|,\ |R_{21}|,\ |R_{22}|
\ll S_2.
\]
Since all norms on $2\times 2$ matrices are equivalent,
\[
\|R_{m,\pm}(s)\|\ll S_2,
\]
uniformly on the microscopic disk. This proves the proposition.
\end{proof}

\begin{proposition}[Entry-by-entry perturbation estimate for the corrected mixed block]
\label{prop:cross-block-perturbation}
Let
\[
q(t)=q_0(t)+\delta q(t),
\qquad
\Phi(t)=\Phi_0(t)+\delta\Phi(t),
\]
where $q_0$ is the baseline local-model-plus-affine-jet phase derivative and
$\delta q$ is the corrected remainder. Write
\[
t_\pm=m\pm \frac{s}{2},
\qquad
\Delta(s):=\Phi(t_-)-\Phi(t_+),
\qquad
\Delta_0(s):=\Phi_0(t_-)-\Phi_0(t_+),
\]
and set
\[
\delta\Delta(s):=\Delta(s)-\Delta_0(s).
\]

Let $N_m^{\corr}(s)$ be the corrected finite-$s$ mixed block built from
$q,\Delta$, and let $N_m^{(0)}(s)$ be the corresponding block built from
$q_0,\Delta_0$. Define
\[
\widetilde R_m(s):=N_m^{\corr}(s)-N_m^{(0)}(s).
\]

Assume that on the microscopic window one has
\[
|\delta q(t_\pm)|\ll L_I^2S_2,
\qquad
|\delta q'(t_\pm)|\ll L_I S_2,
\qquad
|\delta q''(t_\pm)|\ll S_2,
\]
and
\[
|\delta\Delta(s)|\ll L_I^3S_2,
\]
with
\[
L_I\asymp \frac{x_0}{Q},
\qquad
|q_0(t_\pm)|\asymp Q,
\qquad
|s|\asymp L_I.
\]
Then
\[
\boxed{\ \|\widetilde R_m(s)\|\ll S_2\ }
\]
uniformly for $|s|<\rho_0/Q$.
\end{proposition}

\begin{proof}
Write the corrected mixed block entrywise as
\[
N_m^{\corr}(s)=\frac1\pi
\begin{pmatrix}
-\dfrac{2\sin\Delta}{s}
&
\dfrac{\sin\Delta+q(t_+)s\cos\Delta}{s^2}
\\[2ex]
-\dfrac{\sin\Delta+q(t_-)s\cos\Delta}{s^2}
&
\dfrac{(q(t_-)+q(t_+))s\cos\Delta+(2-q(t_-)q(t_+)s^2)\sin\Delta}{2s^3}
\end{pmatrix},
\]
and analogously for $N_m^{(0)}(s)$ with $(q_0,\Delta_0)$.

We estimate the perturbation entry by entry. Throughout, use
\[
\sin(\Delta_0+\delta\Delta)-\sin\Delta_0=O(|\delta\Delta|),
\qquad
\cos(\Delta_0+\delta\Delta)-\cos\Delta_0=O(|\delta\Delta|).
\]

\medskip
\noindent
\textbf{(1,1)-entry.}
\[
\widetilde R_{11}
=
-\frac{2}{\pi s}\bigl(\sin\Delta-\sin\Delta_0\bigr),
\]
hence
\[
|\widetilde R_{11}|
\ll
\frac{|\delta\Delta|}{|s|}
\ll
\frac{L_I^3S_2}{L_I}
\ll
L_I^2S_2
\ll
S_2.
\]

\medskip
\noindent
\textbf{(1,2)-entry.}
The numerator perturbation is
\[
(\sin\Delta-\sin\Delta_0)
+
\bigl(q(t_+)-q_0(t_+)\bigr)s\cos\Delta
+
q_0(t_+)s(\cos\Delta-\cos\Delta_0).
\]
Therefore
\[
|\widetilde R_{12}|
\ll
\frac{|\delta\Delta|+|\delta q(t_+)|\,|s|+|q_0(t_+)|\,|s|\,|\delta\Delta|}{|s|^2}.
\]
Using
\[
|\delta\Delta|\ll L_I^3S_2,\quad
|\delta q(t_+)|\ll L_I^2S_2,\quad
|q_0(t_+)|\asymp Q,\quad
|s|\asymp L_I\asymp Q^{-1},
\]
we get
\[
|\widetilde R_{12}|\ll \frac{L_I^3S_2}{L_I^2}
+\frac{L_I^3S_2}{L_I^2}
+\frac{Q\,L_I^4S_2}{L_I^2}
\ll
\frac{S_2}{Q}
\ll S_2.
\]
The same argument gives
\[
|\widetilde R_{21}|\ll S_2.
\]

\medskip
\noindent
\textbf{(2,2)-entry.}
Write the numerator as
\[
F(q_-,q_+,\Delta;s):=
(q_-+q_+)s\cos\Delta+(2-q_-q_+s^2)\sin\Delta.
\]
Then
\[
\widetilde R_{22}
=
\frac{F(q(t_-),q(t_+),\Delta;s)-F(q_0(t_-),q_0(t_+),\Delta_0;s)}{2\pi s^3}.
\]
A first-order expansion of $F$ in $(q_-,q_+,\Delta)$ gives
\[
|F-F_0|
\ll
|s|\,|\delta q(t_-)|+|s|\,|\delta q(t_+)|
+|s|^2Q\,|\delta q(t_-)|+|s|^2Q\,|\delta q(t_+)|
+\bigl(|s|Q+1\bigr)|\delta\Delta|.
\]
Using
\[
|\delta q(t_\pm)|\ll L_I^2S_2,\qquad
|\delta\Delta|\ll L_I^3S_2,\qquad
|s|\asymp L_I\asymp Q^{-1},
\]
we obtain
\[
|F-F_0|\ll L_I^3S_2.
\]
Hence
\[
|\widetilde R_{22}|
\ll
\frac{L_I^3S_2}{|s|^3}
\ll
S_2.
\]

Collecting the four entrywise bounds and using equivalence of norms on
$2\times2$ matrices gives
\[
\|\widetilde R_m(s)\|\ll S_2,
\]
uniformly on the microscopic disk.
\end{proof}

\begin{lemma}[Baseline same-point variation]
\label{lem:baseline-variation}
Let $G_{m,\pm}^{(0)}(s)$ be the corrected same-point baseline blocks built
from the explicit local model together with the affine jet $J_1[g_{\sm}]$.
Then, on the microscopic disk $|s|<\rho_0/Q$,
\[
\sup_{|s|<\rho_0/Q}
\|G_{m,\pm}^{(0)}(s)-G_{m,\pm}^{(0)}(0)\|
\ll \rho_0 Q.
\]
\end{lemma}

\begin{proof}
Differentiate the explicit same-point block formula entrywise:
\[
G_{m,\pm}^{(0)}(s)=\frac1\pi
\begin{pmatrix}
2q_0(t_\pm) & \frac12 q_0'(t_\pm)\\[1ex]
\frac12 q_0'(t_\pm) & \frac1{12}\bigl(q_0''(t_\pm)+2q_0(t_\pm)^3\bigr)
\end{pmatrix},
\qquad t_\pm=m\pm \frac{s}{2}.
\]
By the chain rule,
\[
\partial_s G_{m,\pm}^{(0)}(s)
=
\pm \frac1{2\pi}
\begin{pmatrix}
2q_0'(t_\pm) & \frac12 q_0''(t_\pm)\\[1ex]
\frac12 q_0''(t_\pm) &
\frac1{12}\bigl(q_0'''(t_\pm)+6q_0(t_\pm)^2q_0'(t_\pm)\bigr)
\end{pmatrix}.
\]
For the explicit local model plus affine jet, one has the baseline local scales
\[
|q_0(t_\pm)|\ll Q,\qquad
|q_0'(t_\pm)|\ll Q^2,\qquad
|q_0''(t_\pm)|\ll Q^3,\qquad
|q_0'''(t_\pm)|\ll Q^4,
\]
uniformly on the admissible midpoint class. The cubic combination
$
q_0'''(t_\pm)+6q_0(t_\pm)^2q_0'(t_\pm)
$
therefore has size $O(Q^2)$ in the current normalization, and hence
\[
\sup_{|s|<\rho_0/Q}\|\partial_s G_{m,\pm}^{(0)}(s)\|\ll Q^2.
\]
Integrating along the segment from $0$ to $s$ gives
\[
\|G_{m,\pm}^{(0)}(s)-G_{m,\pm}^{(0)}(0)\|
\le
|s|\sup_{|u|<\rho_0/Q}\|\partial_u G_{m,\pm}^{(0)}(u)\|
\ll
\frac{\rho_0}{Q}\cdot Q^2
=
\rho_0 Q.
\]
\end{proof}

\begin{proposition}[Corrected finite-$s$ holomorphic whitening]
\label{prop:corrected-whitening}
The corrected same-point and mixed blocks are holomorphic on
\[
|s|<\rho_0/Q.
\]
Moreover, after shrinking $\rho_0$ if necessary, the same-point blocks
\[
G_{m,\pm}^{\corr}(s)
\]
remain uniformly positive definite on this disk, and the inverse square roots
\[
G_{m,\pm}^{\corr}(s)^{-1/2}
\]
are holomorphic there. Consequently
\[
\widehat\Omega_\zeta^{\corr}(s;m)
=
G_{m,-}^{\corr}(s)^{-1/2}
N_m^{\corr}(s)
G_{m,+}^{\corr}(s)^{-1/2}
\]
is holomorphic on $|s|<\rho_0/Q$.
\end{proposition}

\begin{proof}
Holomorphy of the omitted smooth part and its required $t$-derivatives follows
from Proposition~\ref{prop:denominator-comparability}. The only apparent poles
at $s=0$ in the corrected mixed block come from powers of $1/s$, $1/s^2$,
$1/s^3$, and these are removable because $\Delta_m(s)$ is odd in $s$,
$\sin \Delta_m(s)=O(s)$, and the off-diagonal and $(2,2)$ numerators vanish
to orders $2$ and $3$, respectively. Thus both $G_{m,\pm}^{\corr}(s)$ and
$N_m^{\corr}(s)$ are holomorphic on the disk.

By Proposition~\ref{prop:block-perturbation} and the tail-curvature theorem,
\[
\sup_{|s|<\rho_0/Q}\|R_{m,\pm}(s)\|=o(Q).
\]
By Lemma~\ref{lem:baseline-variation},
\[
\sup_{|s|<\rho_0/Q}
\|G_{m,\pm}^{(0)}(s)-G_{m,\pm}^{(0)}(0)\|
\ll \rho_0 Q.
\]
At $s=0$, the baseline jet-limit block has spectral gap
\[
\lambda_{\min}(G_{m,\pm}^{(0)}(0))\gg Q
\]
on the admissible midpoint class. Therefore
\[
\|G_{m,\pm}^{\corr}(s)-G_{m,\pm}^{(0)}(0)\|
\le
\|G_{m,\pm}^{(0)}(s)-G_{m,\pm}^{(0)}(0)\|
+
\|R_{m,\pm}(s)\|
\ll
\rho_0 Q + o(Q).
\]
For $\rho_0$ sufficiently small and $m$ sufficiently large, this is smaller
than half the baseline spectral gap, hence
\[
\lambda_{\min}(G_{m,\pm}^{\corr}(s))\gg Q
\]
uniformly on the disk. Holomorphic functional calculus then gives the
holomorphic inverse square roots, and the whitened block is holomorphic as a
product of holomorphic factors.
\end{proof}

\subsection{Fixed anomaly core and cutoff-dependent auxiliary bookkeeping}
\label{subsec:fixed-core}

To separate the anomaly-carrying local model from the cutoff-dependent tail
bookkeeping, we fix once and for all a finite local core
\[
\mathcal C
\]
containing the candidate off-line anomaly (for example, the designated off-line
pair together with any symmetry-forced companions needed for the local model).

For each cutoff $R\ge R_0$, we then decompose the zero contribution into:
\begin{enumerate}
\item the fixed core contribution $q_{\core}(t)$, coming from the zeros in
$\mathcal C$;
\item the auxiliary near contribution $q_{\aux,R}(t)$, coming from the zeros
outside $\mathcal C$ but inside the near shell determined by $R$;
\item the far contribution $T_{\far,R}(t)$, coming from the zeros outside the
near shell;
\item the background contribution $B_{\bg}(t)$;
\item the estimation term $E_{\est,R}(t)$.
\end{enumerate}

Thus the full phase derivative is decomposed as
\[
q(t)=q_{\core}(t)+q_{\aux,R}(t)+g_{\sm,R}(t)+E_{\est,R}(t),
\]
where
\[
g_{\sm,R}(t):=B_{\bg}(t)+T_{\far,R}(t).
\]

The local toy comparison is attached only to the fixed core $\mathcal C$. The
auxiliary near part $q_{\aux,R}$, the far part $T_{\far,R}$, and the
estimation term $E_{\est,R}$ are absorbed into the corrected zeta-side scalar
defined below.

\begin{proposition}[Corrected local deformation decomposition with fixed core]
\label{prop:corrected-local-decomposition}
Fix the anomaly-carrying core $\mathcal C$ as above, and for each cutoff
$R\ge R_0$ write
\[
q(t)=q_{\core}(t)+q_{\aux,R}(t)+g_{\sm,R}(t)+E_{\est,R}(t),
\qquad
g_{\sm,R}(t)=B_{\bg}(t)+T_{\far,R}(t).
\]
Let $J_1[g_{\sm,R}]$ denote the affine jet of $g_{\sm,R}$ at the midpoint,
and let $\Phi_{\corr,R}$ be the corrected phase obtained from this
decomposition.

Using $\Phi_{\corr,R}$, define the corrected same-point blocks
\[
G_{m,\pm}^{\corr,R}(s),
\]
the corrected mixed block
\[
N_m^{\corr,R}(s),
\]
and the corrected whitened block
\[
\widehat\Omega_\zeta^{\corr,R}(s;m)
=
G_{m,-}^{\corr,R}(s)^{-1/2}
N_m^{\corr,R}(s)
G_{m,+}^{\corr,R}(s)^{-1/2}.
\]

Let
\[
\widehat\Omega_{\bg}^{\corr,R}(s;m)
\]
denote the corresponding corrected whitened block at background value-channel
parameter $S(m)=0$, and define the corrected local deformation
\[
\Delta_\zeta(m,\sigma)
:=
\widehat\Omega_\zeta^{\corr,R}(s;m)-\widehat\Omega_{\bg}^{\corr,R}(s;m),
\qquad s=\sigma.
\]

Then the corrected local deformation admits the expansion
\[
\boxed{
\Delta_\zeta(m,\sigma)
=
S(m)\,A_{\val}(m,\sigma)+R_\zeta(m,\sigma),
}
\]
where $A_{\val}(m,\sigma)$ is the pure value-channel derivative of the
jet-limit block associated with the fixed core $\mathcal C$, and
$R_\zeta(m,\sigma)$ is internal to the full corrected cutoff-dependent scalar
\[
H_{m,R}(s):=\Phi_K\!\bigl(\widehat\Omega_\zeta^{\corr,R}(s;m)\bigr).
\]

In particular, once the corrected odd holomorphic package is established for
$H_{m,R}(s)$, no external tail/error term survives outside the resulting
$N$-point scalar
\[
\Xi_{\zeta,R}^{(N)}(m)
:=
\sum_{j=1}^N a_j^{(N)}H_{m,R}(s_j).
\]
\end{proposition}

\begin{proof}
Consider the analytic whitening map
\[
\mathcal W(G_-,N,G_+):=G_-^{-1/2}NG_+^{-1/2}.
\]
By the corrected finite-$s$ holomorphic-whitening theorem, this map is
holomorphic on the open set where $G_\pm$ remain uniformly positive definite.

For each cutoff $R$, the corrected blocks are built from the decomposition
\[
q=q_{\core}+q_{\aux,R}+g_{\sm,R}+E_{\est,R}.
\]
Thus every auxiliary near contribution from $q_{\aux,R}$, every omitted-zero
tail contribution from $T_{\far,R}$, every affine-jet correction from
$J_1[g_{\sm,R}]$, and every explicit estimation term $E_{\est,R}$ is already
absorbed into the corrected phase $\Phi_{\corr,R}$, hence into the corrected
same-point blocks, the corrected mixed block, the corrected whitened block,
and therefore into the corrected transverse scalar
\[
H_{m,R}(s)=\Phi_K\!\bigl(\widehat\Omega_\zeta^{\corr,R}(s;m)\bigr).
\]

Now expand the corrected whitened block around the background value-channel
parameter $S(m)=0$. By analyticity of $\mathcal W$,
\[
\widehat\Omega_\zeta^{\corr,R}(s;m)
=
\widehat\Omega_{\bg}^{\corr,R}(s;m)
+
S(m)\,D\mathcal W_{(\bg)}[\dot G_-,\dot N,\dot G_+]
+
R_\zeta(m,\sigma),
\]
where the first-order coefficient
\[
D\mathcal W_{(\bg)}[\dot G_-,\dot N,\dot G_+]
\]
is exactly the pure value-channel derivative $A_{\val}(m,\sigma)$ associated
with the fixed core $\mathcal C$, and $R_\zeta(m,\sigma)$ collects the
higher-order terms inside the same corrected cutoff-dependent object.

Subtracting the background block gives
\[
\Delta_\zeta(m,\sigma)
=
S(m)\,A_{\val}(m,\sigma)+R_\zeta(m,\sigma),
\]
as claimed.
\end{proof}

\begin{proposition}[Core-stable cutoff compatibility]
\label{prop:cutoff-compatibility}
Fix the anomaly-carrying core $\mathcal C$, and for each cutoff $R\ge R_0$
define the corrected scalar
\[
H_{m,R}(s):=\Phi_K\!\bigl(\widehat\Omega_\zeta^{\corr,R}(s;m)\bigr)
\]
from the decomposition
\[
q=q_{\core}+q_{\aux,R}+g_{\sm,R}+E_{\est,R}.
\]

Then enlarging $R$ changes only the auxiliary near part $q_{\aux,R}$ and the
far part $T_{\far,R}$, while leaving the anomaly-carrying core
$q_{\core}$ unchanged. Consequently:
\begin{enumerate}
\item the toy anomaly matrix $M(d)$ and the calibration matrix
$A_{\val}(m,\sigma)$ depend only on the fixed core $\mathcal C$, not on
the cutoff $R$;
\item all $R$-dependent auxiliary, tail, jet-correction, and estimation
contributions are internal to the full corrected scalar $H_{m,R}(s)$;
\item if the corrected odd holomorphic package is proved for $H_{m,R}(s)$,
then the $N$-point odd-moment cancellation theorem applies directly to the
full corrected scalar for that cutoff, with no external remainder term.
\end{enumerate}

In particular, the cutoff $R$ may be chosen for quantitative convenience in
the curvature/tail estimates without changing the fixed toy comparison
attached to the core $\mathcal C$.
\end{proposition}

\begin{proof}
Let $R_1<R_2$. Then the only change in the decomposition
\[
q=q_{\core}+q_{\aux,R}+g_{\sm,R}+E_{\est,R}
\]
is that some non-core zeros move from the far part into the auxiliary near
part:
\[
q_{\aux,R_2}-q_{\aux,R_1}
=
-\bigl(T_{\far,R_2}-T_{\far,R_1}\bigr).
\]
The core contribution $q_{\core}$ is unchanged by construction.

Thus the toy anomaly matrix $M(d)$, which is attached only to the fixed core,
is unchanged, and the value-channel derivative $A_{\val}(m,\sigma)$ associated
with that core is unchanged as well. All cutoff dependence is confined to the
auxiliary near contribution, the far tail, the affine-jet correction built
from $g_{\sm,R}$, and the explicit estimation term $E_{\est,R}$. These enter
only through the corrected phase $\Phi_{\corr,R}$, hence only through the full
corrected scalar $H_{m,R}(s)$.

Therefore the corrected odd holomorphic package, once established for
$H_{m,R}(s)$, applies to the complete cutoff-dependent object. The resulting
$N$-point scalar already contains all auxiliary/tail/error effects for that
cutoff, and no external remainder survives. Since the fixed core is
unchanged, the toy comparison is stable under enlarging $R$. This proves the
claim.
\end{proof}

\begin{remark}[Cutoff choice]
\label{rem:cutoff-choice}
Proposition~\ref{prop:cutoff-compatibility} is a structural statement. It
does not assert any particular optimal choice of $R$; it only says that once
the core $\mathcal C$ is fixed, the cutoff may be adjusted in the perturbative
bookkeeping without altering the underlying toy comparison.
\end{remark}

\begin{proposition}[Whitened mixed-block transfer with preserved $Q^{-3}$ scale]
\label{prop:whitened-mixed-transfer}
Let
\[
\mathcal W(G_-,N,G_+):=G_-^{-1/2}NG_+^{-1/2}.
\]
Write the corrected finite-$s$ blocks as
\[
G_\pm^{\corr}(s)=G_\pm^{(0)}(s)+R_\pm(s),
\qquad
N^{\corr}(s)=N^{(0)}(s)+\widetilde R(s),
\]
and set
\[
\widehat\Omega_\zeta^{\corr}(s;m):=\mathcal W(G_-^{\corr},N^{\corr},G_+^{\corr}),
\qquad
\widehat\Omega_\zeta^{(0)}(s;m):=\mathcal W(G_-^{(0)},N^{(0)},G_+^{(0)}).
\]

Then
\[
\widehat\Omega_\zeta^{\corr}(s;m)-\widehat\Omega_\zeta^{(0)}(s;m)
=
\mathcal T_-(s;m)+\mathcal T_{\mix}(s;m)+\mathcal T_+(s;m)+\mathcal E_2(s;m),
\]
where
\[
\mathcal T_-(s;m):=X_-(s)N^{(0)}(s)G_+^{(0)}(s)^{-1/2},
\]
\[
\mathcal T_{\mix}(s;m):=G_-^{(0)}(s)^{-1/2}\widetilde R(s)G_+^{(0)}(s)^{-1/2},
\]
\[
\mathcal T_+(s;m):=G_-^{(0)}(s)^{-1/2}N^{(0)}(s)X_+(s),
\]
with
\[
X_\pm(s):=D(G^{-1/2})_{G_\pm^{(0)}(s)}[R_\pm(s)],
\]
and $\mathcal E_2$ quadratic in $(R_-,\widetilde R,R_+)$.

Moreover,
\[
\Phi_K(\mathcal T_-(s;m))\ll \frac{S_2}{Q^3},
\qquad
\Phi_K(\mathcal T_{\mix}(s;m))\ll \frac{S_2}{Q^3},
\qquad
\Phi_K(\mathcal T_+(s;m))\ll \frac{S_2}{Q^3},
\]
and
\[
\Phi_K(\mathcal E_2(s;m))\ll \frac{S_2^2}{Q^3}.
\]
Hence
\[
\boxed{
\Phi_K\!\bigl(\widehat\Omega_\zeta^{\corr}(s;m)-\widehat\Omega_\zeta^{(0)}(s;m)\bigr)
\ll
\frac{S_2}{Q^3}.
}
\]
\end{proposition}

\begin{proof}
We first collect the scales used below from results established earlier. The
same-point perturbation bound $\|R_\pm(s)\|\ll S_2$ is
Proposition~\ref{prop:block-perturbation}; the mixed-block perturbation bound
$\|\widetilde R(s)\|\ll S_2$ is
Proposition~\ref{prop:cross-block-perturbation}; and the smallness
$S_2=o(Q)$ is Theorem~\ref{thm:tail-curvature}. The baseline entry scales
\[
G_\pm^{(0)}(s)=
\begin{pmatrix} O(Q) & O(1)\\ O(1) & O(Q^3) \end{pmatrix},
\qquad
G_\pm^{(0)}(s)^{-1/2}=
\begin{pmatrix}
O(Q^{-1/2}) & O(Q^{-3/2})\\
O(Q^{-3/2}) & O(Q^{-3/2})
\end{pmatrix},
\]
\[
N^{(0)}(s)=
\begin{pmatrix} O(Q) & O(Q^2)\\ O(Q^2) & O(Q^3) \end{pmatrix}
\]
are read off directly from the explicit same-point formula \eqref{eq:Gpm} and
the mixed-block formula \eqref{eq:Nm}, given the baseline local scales of
$q_0, q_0', q_0''$ stated in the proof of
Lemma~\ref{lem:baseline-variation}. Uniform positive-definiteness of
$G_\pm^{(0)}(s)$ on the microscopic disk is
Proposition~\ref{prop:corrected-whitening}.

The whitening map admits the first-order expansion
\[
\widehat\Omega_\zeta^{\corr}-\widehat\Omega_\zeta^{(0)}
=
D\mathcal W_{(0)}[R_-,\widetilde R,R_+]+\mathcal E_2,
\]
where
\[
D\mathcal W_{(0)}[R_-,\widetilde R,R_+]
=
X_-N^{(0)}G_+^{(0)-1/2}
+
G_-^{(0)-1/2}\widetilde R\,G_+^{(0)-1/2}
+
G_-^{(0)-1/2}N^{(0)}X_+.
\]

We estimate the three first-order pieces separately.

\medskip
\noindent
\textbf{Step 1: the mixed-block term.}
By Proposition~\ref{prop:cross-block-perturbation} (entrywise form),
\[
\widetilde R_{11}\ll \frac{S_2}{Q^2},
\qquad
\widetilde R_{12},\widetilde R_{21}\ll \frac{S_2}{Q},
\qquad
\widetilde R_{22}\ll S_2.
\]
Combining with the baseline inverse-square-root scales above, the off-diagonal
entries of
\[
\mathcal T_{\mix}:=G_-^{(0)-1/2}\widetilde R\,G_+^{(0)-1/2}
\]
are all $O(S_2/Q^3)$. Hence
\[
\Phi_K(\mathcal T_{\mix})\ll \frac{S_2}{Q^3}.
\]

\medskip
\noindent
\textbf{Step 2: the same-point Fréchet derivative.}
We prove the corresponding estimate for the left same-point term
\[
\mathcal T_-:=X_-N^{(0)}G_+^{(0)-1/2},
\qquad
X_-:=D(G^{-1/2})_{G_-^{(0)}}[R_-].
\]
The proof for $\mathcal T_+$ is identical.

Write one baseline same-point block in the form
\[
G_0=
\begin{pmatrix}
a & b\\
b & c
\end{pmatrix},
\qquad
a\asymp Q,\quad b=O(1),\quad c\asymp Q^3,
\]
and one perturbation as
\[
R=
\begin{pmatrix}
r_{11} & r_{12}\\
r_{12} & r_{22}
\end{pmatrix},
\qquad
r_{11}\ll \frac{S_2}{Q^2},\quad
r_{12}\ll \frac{S_2}{Q},\quad
r_{22}\ll S_2.
\]

Diagonalize
\[
G_0=U\Lambda U^\top,
\qquad
\Lambda=\operatorname{diag}(\lambda_1,\lambda_2),
\]
with $U$ orthogonal. Since
\[
a\asymp Q,\qquad c\asymp Q^3,\qquad b=O(1),
\]
we have
\[
\lambda_1\asymp Q,\qquad \lambda_2\asymp Q^3.
\]
Moreover, if $U$ is a rotation by angle $\theta$, then
\[
\tan(2\theta)=\frac{2b}{c-a}=O(Q^{-3}),
\]
hence
\[
\theta=O(Q^{-3}),
\qquad
U=I+O(Q^{-3}).
\]

Set
\[
H:=U^\top R\,U.
\]
A direct computation gives
\[
H_{11}=r_{11}+2\theta r_{12}+O(\theta^2 r_{22}),
\]
\[
H_{12}=r_{12}+\theta(r_{22}-r_{11})+O(\theta^2 r_{12}),
\]
\[
H_{22}=r_{22}-2\theta r_{12}+O(\theta^2 r_{11}).
\]
Using
\[
\theta=O(Q^{-3}),
\qquad
r_{11}\ll \frac{S_2}{Q^2},
\qquad
r_{12}\ll \frac{S_2}{Q},
\qquad
r_{22}\ll S_2,
\]
we obtain
\[
H_{11}\ll \frac{S_2}{Q^2},
\qquad
H_{12}\ll \frac{S_2}{Q},
\qquad
H_{22}\ll S_2.
\]

Now let $f(x)=x^{-1/2}$. For a diagonal matrix $\Lambda$, the Fréchet
derivative is given entrywise by the Löwner formula
\[
Df_\Lambda(H)_{ij}=f^{[1]}(\lambda_i,\lambda_j)H_{ij},
\]
where
\[
f^{[1]}(u,v)=
\begin{cases}
\dfrac{u^{-1/2}-v^{-1/2}}{u-v}, & u\neq v,\\[1.2ex]
-\dfrac{1}{2u^{3/2}}, & u=v.
\end{cases}
\]
Since
\[
\lambda_1\asymp Q,\qquad \lambda_2\asymp Q^3,
\]
we have
\[
f^{[1]}(\lambda_1,\lambda_1)\asymp Q^{-3/2},
\qquad
f^{[1]}(\lambda_2,\lambda_2)\asymp Q^{-9/2},
\qquad
f^{[1]}(\lambda_1,\lambda_2)\asymp Q^{-7/2}.
\]
Therefore, if
\[
Y:=Df_\Lambda(H),
\]
then
\[
Y_{11}\ll Q^{-3/2}\cdot \frac{S_2}{Q^2}
=\frac{S_2}{Q^{7/2}},
\]
\[
Y_{12}\ll Q^{-7/2}\cdot \frac{S_2}{Q}
=\frac{S_2}{Q^{9/2}},
\]
\[
Y_{22}\ll Q^{-9/2}\cdot S_2
=\frac{S_2}{Q^{9/2}}.
\]

Since
\[
X=U\,Y\,U^\top
\]
and $U=I+O(Q^{-3})$, these scales persist in the original basis:
\[
X_{11}\ll \frac{S_2}{Q^{7/2}},
\qquad
X_{12},X_{21}\ll \frac{S_2}{Q^{9/2}},
\qquad
X_{22}\ll \frac{S_2}{Q^{9/2}}.
\]

\medskip
\noindent
\textbf{Step 3: the left same-point term $\mathcal T_-$.}
The baseline mixed-block scales are
\[
N^{(0)}=
\begin{pmatrix}
O(Q) & O(Q^2)\\
O(Q^2) & O(Q^3)
\end{pmatrix},
\qquad
G_+^{(0)-1/2}=
\begin{pmatrix}
O(Q^{-1/2}) & O(Q^{-3/2})\\
O(Q^{-3/2}) & O(Q^{-3/2})
\end{pmatrix}.
\]
We estimate only the off-diagonal entries, since
\[
\Phi_K(Y)=Y_{12}+Y_{21}.
\]

For $(\mathcal T_-)_{12}$, every summand is $O(S_2/Q^3)$. For example,
\[
X_{11}N_{12}(G_+^{-1/2})_{22}
\ll
\frac{S_2}{Q^{7/2}}\cdot Q^2\cdot Q^{-3/2}
=
\frac{S_2}{Q^3},
\]
and
\[
X_{12}N_{22}(G_+^{-1/2})_{22}
\ll
\frac{S_2}{Q^{9/2}}\cdot Q^3\cdot Q^{-3/2}
=
\frac{S_2}{Q^3}.
\]
All other terms are smaller. Thus
\[
(\mathcal T_-)_{12}\ll \frac{S_2}{Q^3}.
\]
The same argument gives
\[
(\mathcal T_-)_{21}\ll \frac{S_2}{Q^3},
\]
and hence
\[
\Phi_K(\mathcal T_-)\ll \frac{S_2}{Q^3}.
\]

By symmetry, the same estimate holds for the right same-point term:
\[
\Phi_K(\mathcal T_+)\ll \frac{S_2}{Q^3}.
\]

\medskip
\noindent
\textbf{Step 4: the quadratic remainder.}
Every term in $\mathcal E_2$ contains at least two perturbation factors from
$(R_-,\widetilde R,R_+)$, and whitening contributes only polynomial
$Q$-weights. Therefore
\[
\Phi_K(\mathcal E_2)\ll \frac{S_2^2}{Q^3}.
\]
By Theorem~\ref{thm:tail-curvature}, $S_2=o(Q)$, so this is lower order than
$S_2/Q^3$.

Combining the estimates for $\mathcal T_-$, $\mathcal T_{\mix}$,
$\mathcal T_+$, and $\mathcal E_2$ gives
\[
\Phi_K\!\bigl(\widehat\Omega_\zeta^{\corr}(s;m)-\widehat\Omega_\zeta^{(0)}(s;m)\bigr)
\ll
\frac{S_2}{Q^3},
\]
as claimed.
\end{proof}

\section{Odd holomorphic transverse channel}

Define the corrected transverse scalar
\begin{equation}\label{eq:Hm-def}
H_m(s):=\Phi_K\!\bigl(\widehat\Omega_\zeta^{\corr}(s;m)\bigr).
\end{equation}
By symmetry under $s\mapsto -s$, this is odd. By
Corollary~\ref{cor:PhiK-Aval-zero}, the leading value channel is
annihilated by $\Phi_K$, so only the curvature, tail, and corrected
finite-$s$ remainder channels contribute.

\begin{proposition}[Boundary estimate]
\label{prop:boundary-estimate}
On the microscopic boundary
\[
|s|=\rho_0/Q,
\]
one has
\[
|H_m(s)|\ll_{\rho_0}\frac{L(m)S(m)^2}{Q^3}.
\]
\end{proposition}

\begin{proof}
By Proposition~\ref{prop:corrected-local-decomposition}, the corrected local
deformation may be written in the form
\[
\Delta_\zeta(m,\sigma)
=
S(m)\,A_{\val}(m,\sigma)+R_\zeta(m,\sigma),
\]
where $R_\zeta(m,\sigma)$ is internal to the full corrected cutoff-dependent
scalar
\[
H_{m,R}(s):=\Phi_K\!\bigl(\widehat\Omega_\zeta^{\corr,R}(s;m)\bigr).
\]
By Proposition~\ref{prop:cutoff-compatibility}, all cutoff-dependent tail,
jet-correction, and estimation contributions are already absorbed into this
full corrected scalar, so no external remainder term is carried separately at
the boundary stage.

Applying $\Phi_K$ to the corrected local decomposition and using
Corollary~\ref{cor:PhiK-Aval-zero},
\[
\Phi_K(A_{\val}(m,\sigma))=0,
\]
gives
\[
H_{m,R}(s)
=
\Phi_K\!\bigl(\widehat\Omega_\zeta^{\corr,R}(s;m)\bigr)
=
\Phi_K\!\bigl(R_\zeta(m,\sigma)\bigr).
\]

Now compare the full corrected whitened block with its baseline counterpart.
By Proposition~\ref{prop:whitened-mixed-transfer},
\[
\Phi_K\!\bigl(\widehat\Omega_\zeta^{\corr,R}(s;m)-\widehat\Omega_\zeta^{(0),R}(s;m)\bigr)
\ll
\frac{S_2}{Q^3}
\qquad (|s|<\rho_0/Q).
\]
On the microscopic boundary $|s|=\rho_0/Q$, the tail-curvature theorem and
the corrected transfer bounds imply
\[
S_2\ll L(m)S(m)^2
\]
in the corrected regime. Therefore
\[
|H_{m,R}(s)|
\ll
\frac{L(m)S(m)^2}{Q^3}
\qquad (|s|=\rho_0/Q).
\]

Since the cutoff bookkeeping has already been absorbed into the full corrected
scalar $H_{m,R}(s)$, this is the boundary estimate for the complete corrected
object, not merely for a leading part plus an external tail/error remainder.
This proves the claim.
\end{proof}

\begin{proposition}[Odd microscopic expansion]
\label{prop:odd-expansion}
The function $H_m(s)$ is odd and holomorphic on $|s|<\rho_0/Q$, and therefore
admits an expansion
\[
H_m(z/Q^2)
=
\sum_{r\ge 0} c_{2r+1}(m)\frac{z^{2r+1}}{Q^{2r+4}},
\qquad |z|<\rho_0,
\]
with
\[
|c_{2r+1}(m)|\ll_{\rho_0}L(m)S(m)^2\,\rho_0^{-(2r+1)}.
\]
\end{proposition}

\begin{proof}
Oddness follows from symmetric placement and the jet-basis swap symmetry:
\[
\widehat\Omega_\zeta^{\corr}(-s;m)=J_0\,\widehat\Omega_\zeta^{\corr}(s;m)\,J_0,
\]
with
\[
\Phi_K(J_0XJ_0)=-\Phi_K(X).
\]
Holomorphy follows from Proposition~\ref{prop:corrected-whitening}, and the
coefficient bounds are immediate from Cauchy's estimate applied to
Proposition~\ref{prop:boundary-estimate}.
\end{proof}

\section{\texorpdfstring{$N$}{N}-point odd-moment cancellation}

For $N\le \rho_0Q$, set
\[
s_j=\frac{j}{Q^2},
\qquad j=1,\dots,N,
\]
and let $a_j^{(N)}$ be the explicit odd-moment-canceling coefficients satisfying
\[
\sum_{j=1}^N a_j^{(N)}=1,
\qquad
\sum_{j=1}^N a_j^{(N)}j^{2r+1}=0
\quad (0\le r\le N-2).
\]
Define the corrected $N$-point transverse scalar
\begin{equation}\label{eq:Xi-zeta}
\Xi_\zeta^{(N)}(m):=
\sum_{j=1}^N a_j^{(N)}H_m(s_j).
\end{equation}

\begin{proposition}[Corrected $N$-point transverse estimate for the full corrected scalar]
\label{prop:corrected-n-point}
Let
\[
H_m(s):=\Phi_K\!\bigl(\widehat\Omega_\zeta^{\corr}(s;m)\bigr)
\]
be the full corrected transverse scalar.
For $N=\lfloor \alpha Q\rfloor$ with $\alpha>0$ sufficiently small,
define
\[
s_j=\frac{j}{Q^2},
\qquad
\Xi_\zeta^{(N)}(m)
:=
\sum_{j=1}^N a_j^{(N)}\,H_m(s_j).
\]
Then
\[
|\Xi_\zeta^{(N)}(m)|\ll L(m)S(m)^2e^{-\delta Q}
\]
for some $\delta>0$.
\end{proposition}

\begin{proof}
By Proposition~\ref{prop:corrected-whitening}, the full corrected scalar
\[
H_m(s)=\Phi_K\!\bigl(\widehat\Omega_\zeta^{\corr}(s;m)\bigr)
\]
is holomorphic on $|s|<\rho_0/Q$. By
Proposition~\ref{prop:cutoff-compatibility}, all cutoff-dependent tail,
jet-correction, and estimation contributions are already internal to the full
corrected scalar $H_m(s)$; hence no external remainder is carried into
$\Xi_\zeta^{(N)}(m)$. By Proposition~\ref{prop:boundary-estimate},
it satisfies
\[
|H_m(s)|\ll_{\rho_0}\frac{L(m)S(m)^2}{Q^3}
\qquad (|s|=\rho_0/Q).
\]
By Proposition~\ref{prop:odd-expansion}, it is odd and admits the expansion
\[
H_m(z/Q^2)
=
\sum_{r\ge 0} c_{2r+1}(m)\frac{z^{2r+1}}{Q^{2r+4}},
\]
with
\[
|c_{2r+1}(m)|\ll L(m)S(m)^2\,\rho_0^{-(2r+1)}.
\]

Applying the standard odd-moment cancellation with the coefficients
$a_j^{(N)}$ kills the first $N-1$ odd terms of this full corrected
expansion. The surviving moment factor is exponentially tame for
$N=\alpha Q$, which gives
\[
|\Xi_\zeta^{(N)}(m)|\ll L(m)S(m)^2e^{-\delta Q}.
\]
\end{proof}

We isolate the toy microscopic expansion in
Appendix~\ref{app:toy-microscopic-expansion}; the only fact used in the main
argument is the following $N$-point preservation statement.

\begin{proposition}[Toy $N$-point transverse preservation]
\label{prop:toy-n-point-direct}
Fix a compact interval
\[
D=[d_-,d_+]\subset(0,\infty),
\]
and let
\[
\Xi_{\toy}^{(N)}(u,d)
:=
\sum_{j=1}^N a_j^{(N)}\,
\Phi_K\!\bigl(\Delta_{\toy}(u,d;s_j)\bigr),
\qquad
s_j=\frac{j}{Q^2},
\]
where the coefficients satisfy
\[
\sum_{j=1}^N a_j^{(N)}=1.
\]
Then, for $d\in D$ and $|u|<u_0(D)$,
\[
\Xi_{\toy}^{(N)}(u,d)
=
u^2\Phi_K(M(d))+O(u^4),
\]
uniformly in $d\in D$. In particular,
\[
|\Xi_{\toy}^{(N)}(u,d)|\asymp u^2
\]
uniformly on compact $d$-ranges.
\end{proposition}

\begin{proof}
By Lemma~\ref{lem:toy-microscopic-expansion}, one has
\[
\Delta_{\toy}(u,d;s_j)=u^2M(d)+u^4E_j(u,d),
\qquad
\sup_{1\le j\le N}\|E_j(u,d)\|\ll 1
\]
uniformly for $d\in D$. Applying $\Phi_K$ gives
\[
\Phi_K\!\bigl(\Delta_{\toy}(u,d;s_j)\bigr)
=
u^2\Phi_K(M(d))+u^4\Phi_K(E_j(u,d)).
\]
Summing over $j$ and using
\[
\sum_{j=1}^N a_j^{(N)}=1
\]
yields
\[
\Xi_{\toy}^{(N)}(u,d)
=
u^2\Phi_K(M(d))
+
u^4\sum_{j=1}^N a_j^{(N)}\Phi_K(E_j(u,d)).
\]
Since $\Phi_K$ is bounded and $\sup_j\|E_j(u,d)\|\ll 1$, the error term is
$O(u^4)$, uniformly for $d\in D$. This proves the expansion.

Because $\Phi_K(M(d))\neq 0$ for $d>0$, and is bounded away from $0$ on every
compact interval $D\subset(0,\infty)$, the final estimate
\[
|\Xi_{\toy}^{(N)}(u,d)|\asymp u^2
\]
follows.
\end{proof}

On the toy side, Proposition~\ref{prop:toy-n-point-direct} gives
\begin{equation}\label{eq:Xi-toy}
\Xi_{\toy}^{(N)}(u,d)=u^2\Phi_K(M(d))+O(u^4),
\end{equation}
so on compact $d$-ranges,
\[
|\Xi_{\toy}^{(N)}(u,d)|\asymp u^2.
\]
Using \eqref{eq:u2-calibration},
\begin{equation}\label{eq:Xi-toy-calibrated}
|\Xi_{\toy}^{(N)}(u,d)|\asymp \frac{x_0}{B_\zeta(m)}S(m).
\end{equation}

\section{Final contradiction}

The calibrated $N$-point comparison now closes the argument.

\begin{lemma}[Crude polynomial bound for the local value scale]
\label{lem:crude-S-bound}
Let
\[
S(m)=\sum_\rho f_{\beta_\rho,\gamma_\rho}(m),
\]
where
\[
f_{\beta,\gamma}(m)
=
\frac{1-\beta}{(1-\beta)^2+(2m-\gamma)^2}
+
\frac{\beta}{\beta^2+(2m-\gamma)^2}.
\]
Then
\[
S(m)\ll Q^2,
\qquad Q\asymp \log|m|.
\]
Consequently
\[
L(m)\ll Q
\]
and
\[
B_\zeta(m)L(m)S(m)\ll Q^4.
\]
\end{lemma}

\begin{proof}
By the zero-free region, for zeros at relevant height one has
\[
\min(\beta,1-\beta)\gg \frac1Q.
\]
Hence
\[
f_{\beta,\gamma}(m)\ll \frac{Q}{1+(2m-\gamma)^2}.
\]
Partition the zeros into unit shells
\[
\mathcal S_k:=\{\rho: k\le |2m-\gamma_\rho|<k+1\},
\qquad k=0,1,2,\dots
\]
Standard zero counting gives
\[
\#\mathcal S_k\ll Q
\]
for each unit shell near height $m$. Therefore
\[
S(m)
\ll
\sum_{k\ge 0}\#\mathcal S_k\frac{Q}{1+k^2}
\ll
Q\sum_{k\ge 0}\frac{Q}{1+k^2}
\ll
Q^2.
\]
Since
\[
L(m)=\log\!\Bigl(3+|m|+\frac1{S(m)}\Bigr),
\]
we trivially have $L(m)\ll Q$, and the last estimate follows from
$B_\zeta(m)\asymp Q$.
\end{proof}

\begin{theorem}[Final calibrated $N$-point contradiction]
\label{thm:final-contradiction}
No fixed off-line toy anomaly can occur. Consequently the Riemann Hypothesis
holds.
\end{theorem}

\begin{proof}
By Proposition~\ref{prop:calibration-remainder-cutoff} and the canonical
calibration theorem, in the fixed small-$x_0$ regime one has
\[
u^2\asymp \frac{x_0}{B_\zeta(m)}S(m).
\]

On the toy side, Proposition~\ref{prop:toy-n-point-direct} gives
\[
\Xi_{\toy}^{(N)}(u,d)=u^2\Phi_K(M(d))+O(u^4),
\]
hence on compact $d$-ranges,
\[
|\Xi_{\toy}^{(N)}(u,d)|\asymp u^2
\asymp
\frac{x_0}{B_\zeta(m)}S(m).
\]

On the zeta side, Proposition~\ref{prop:corrected-n-point} gives
\[
|\Xi_\zeta^{(N)}(m)|\ll L(m)S(m)^2e^{-\delta Q}.
\]

By Lemma~\ref{lem:crude-S-bound}, one has $S(m)\ll Q^2$, $L(m)\ll Q$, and
$B_\zeta(m)L(m)S(m)\ll Q^4$. Therefore
\[
\frac{L(m)S(m)^2e^{-\delta Q}}{(x_0/B_\zeta(m))S(m)}
=
\frac{B_\zeta(m)}{x_0}L(m)S(m)e^{-\delta Q}
\ll
Q^4e^{-\delta Q}\to 0.
\]
Hence for $Q$ sufficiently large,
\[
|\Xi_\zeta^{(N)}(m)|<|\Xi_{\toy}^{(N)}(u,d)|.
\]

This is impossible if the corrected local zeta deformation realizes the local
off-line toy anomaly with parameter $d\in D$, because the calibrated $N$-point
construction was set up so that the toy model is the jet-limit local avatar of
an off-line zero configuration. Thus no off-line zero can occur. Since every
violation of the Riemann Hypothesis would produce such an off-line local
anomaly, all nontrivial zeros must lie on the critical line. Hence $\RH$
holds.
\end{proof}

\appendix

\section{Toy microscopic expansion}
\label{app:toy-microscopic-expansion}

\begin{lemma}[Uniform microscopic expansion of the toy local deformation]
\label{lem:toy-microscopic-expansion}
Fix a compact interval
\[
D=[d_-,d_+]\subset(0,\infty),
\]
and let
\[
\Delta_{\toy}(u,d;s)
\]
denote the toy local deformation matrix in the symmetric microscopic regime,
with off-line parameter $u$, separation parameter $d$, and microscopic node
$s$.

Then there exist constants $\rho_{\toy}>0$ and $u_0(D)>0$ such that for all
\[
d\in D,
\qquad
|s|<\rho_{\toy},
\qquad
|u|<u_0(D),
\]
one has the expansion
\[
\Delta_{\toy}(u,d;s)=u^2M(d)+u^4E(u,d;s),
\]
where $M(d)$ is the toy anomaly matrix and
\[
\sup_{\substack{d\in D\\ |s|<\rho_{\toy}\\ |u|<u_0(D)}} \|E(u,d;s)\|\ll 1.
\]
In particular, for the microscopic nodes
\[
s_j=\frac{j}{Q^2},
\qquad
j=1,\dots,N,
\]
with $N\le \rho_{\toy}Q^2$, one has uniformly
\[
\Delta_{\toy}(u,d;s_j)=u^2M(d)+u^4E_j(u,d),
\qquad
\sup_j \|E_j(u,d)\|\ll 1.
\]
\end{lemma}

\begin{proof}
The toy local deformation is obtained from an explicit finite-dimensional
off-line perturbation of the toy block. By construction it depends smoothly
(in fact analytically) on the parameters $(u,s,d)$ in a neighborhood of
\[
(u,s,d)=(0,0,d_0),
\qquad d_0\in D.
\]
At $u=0$, the deformation vanishes, and by the even symmetry of the off-line
split there is no linear term in $u$. Hence the first nontrivial term is
quadratic:
\[
\Delta_{\toy}(u,d;s)=u^2M(d,s)+O(u^4),
\]
uniformly on compact $d$-ranges and for $|s|$ in a fixed microscopic disk.

By the jet-limit computation for the toy model, the leading quadratic
coefficient is independent of the microscopic node $s$ and equals the toy
anomaly matrix $M(d)$. Thus
\[
M(d,s)\equiv M(d),
\]
and the expansion becomes
\[
\Delta_{\toy}(u,d;s)=u^2M(d)+u^4E(u,d;s),
\]
with $E$ uniformly bounded on compact parameter ranges. Restricting to the
microscopic nodes $s_j=j/Q^2$ gives the stated uniform bound for
$E_j(u,d)$.
\end{proof}

\end{document}
